\begin{textsample}{NEG Dim 1 – 2023\_c5 – Score 21.00 – 2023\_c5/enade\_2023\_eng\_qui\_7565c19e5a.txt}  \label{ex:f1_neg_429}
CAPÍTULO 4 DISTRIBUIÇÃO DE CONCEITOS A finalidade deste capítulo é apresentar um panorama nacional sobre a distribuição do Conceito Enade 2023.
Para traçar esse panorama, recorre -se a tabelas e gráficos, que são acompanhados de comentários e análises que contribuem para a contextualização dos dados neles expostos.
Para visualizar o panorama do Enade 2023, em nível nacional, na primeira seção deste capítulo, os dados sobre o número de cursos participantes são organizados por grande região.
Na segunda, o mapeamento do conceito do exame é feito por categoria administrativa das IES e modalidade de oferta dos cursos, de maneira cruzada às grandes regiões.
Na última seção, a lógica de apresentação dos dados é feita por organização acadêmica e, também, por grande região de funcionamento dos cursos, no país como um todo. 4.1 PANORAMA NACIONAL DA DISTRIBUIÇÃO DOS CONCEITOS A lógica adotada na apresentação dos dados, nesta seção e nas duas seguintes, pressupõe o registro de duas informações de natureza reiterativa.
O primeiro diz respeito aos valores possíveis do Enade, que variam em níveis de 1 a 5, sendo que, à medida que esses valores aumentam, o curso é melhor avaliado, e ao fato de que aqueles cursos que tiveram somente um ou nenhum concluinte participante não obtêm o Conceito Enade, ficando sem conceito ( SC ), como ressaltado na apresentação deste relatório.
O segundo registro é que, caso haja diferença entre os cursos tabulados neste capítulo e no capítulo 1, ela corresponderá aos cursos que não foram avaliados, ou seja, àqueles, em princípio, sem estudantes concluintes que tivessem sido inscritos para a prova.
Na Tabela 4.1, são apresentadas a quantidade e a distribuição de cursos de Engenharia Química participantes do Enade 2023, por faixa de conceito e grande região.
A partir dos dados da Tabela 4.1, nota -se que, dos 185 cursos participantes, 70 ( 37,8 \% ) classificaram -se com conceito 3, o valor modal.
Esse foi também o conceito modal em todas nas regiões : Nordeste ( 31,0 \% ), Norte ( 71,4 \% ), Sudeste ( 35,6 \% ), Sul ( 42,9 \% ).
Para a região Centro-Oeste, o conceito modal foi 4 ( 33,3 \% ).
O conceito 2 foi o segundo mais frequente em nível nacional ( 28,6 \%, correspondendo a 53 cursos ), tendo sido o segundo conceito mais frequente na maioria das regiões.
O conceito 4 foi o terceiro ( 15,7 \%, correspondendo a 29 cursos ).
Houve, ainda, 9 ( 4,9 \% ) cursos que obtiveram o conceito 5 e oito ( 4,3 \% ) cursos que obtiveram o conceito 1.
Dos 185 cursos de Engenharia Química, 16 ( 8,6 \% ) ficaram sem conceito ( SC ).
A região Norte participou com 7 cursos, o que corresponde a 3,8 \% do total nacional.
Desses, cinco cursos alcançaram o conceito 3, o que equivale a 71,4 \% do total regional, o conceito modal na região, conforme mencionado.
A um curso ( 14,3 \% ) atribuiu -se o conceito 2 ; a um curso ( 14,3 \% ), o conceito 1 ; e a nenhum curso ( 0,0 \% ) o conceito 5.
Ainda nessa região, nenhum curso ( 0,0 \% ) obteve conceito 4 e nenhum curso ( 0,0 \% ) \textbf{esteve} sem conceito ( SC ).
A região Nordeste participou com 29 cursos, correspondentes a 15,7 \% do total nacional.
Nessa região, nove cursos ( 31,0 \% em termos regionais ) obtiveram conceito 3, o conceito modal para a região, como já citado.
Os conceitos 2 e 4 foram atribuídos, respectivamente, a sete e a seis cursos ( 24,1 \% e 20,7 \%, respectivamente ).
Ainda nessa região, um curso ( 3,4 \% ) foi avaliado com o conceito 1 e nenhum curso ( 0,0 \% ) foi avaliado com o conceito 5.
Outros seis cursos ( 20,7 \% ) ficaram sem conceito ( SC ).
Para os 87 cursos participantes da região Sudeste ( 47,0 \% do total nacional ), o conceito 3 foi o conceito modal, de forma já mencionada anteriormente, tendo sido obtido por 31 cursos ( 35,6 \% ).
O conceito 2 foi atribuído a 28 cursos ( 32,2 \% ) e o conceito 4, a 13 cursos ( 14,9 \% ).
Receberam os conceitos 5 e 1, respectivamente, quatro cursos ( 4,6 \% ) e três cursos ( 3,4 \% ).
Nessa região, oito cursos ( 9,2 \% ) ficaram sem conceito ( SC ).
Os 56 cursos da região Sul corresponderam a 30,3 \% do total nacional.
Como já indicado, houve predominância do conceito 3 ( conceito modal ), atribuído a 42,9 \% dos cursos da região, ou seja, atribuído a 24 dos 56 cursos participantes na região Sul.
O conceito 2 foi atribuído a 16 cursos ( 28,6 \% ) ; o conceito 4, oito cursos ( 14,3 \% ) ; e o conceito 5, a quatro cursos ( 7,1 \% ).
Na região Sul, dois cursos ( 3,6 \% ) obtiveram conceito 1 e dois cursos ( 3,6 \% ) ficaram sem conceito ( SC ).
Já dos 6 cursos participantes na região Centro-Oeste ( 3,2 \% do total nacional ), 2 ( 33,3 \% em termos regionais ) receberam conceito 4, o conceito modal, como já destacado.
Um curso ( 16,7 \% ) obteve o conceito 5, um curso ( 16,7 \% ) obtive o conceito 3 e um curso ( 16,7 \% ), o conceito 2.
Nessa região, um curso ( 16,7 \% ) obteve o conceito 1 e nenhum curso ( 0,0 \% ) ficou sem conceito ( SC ).
Em um gráfico de distribuição cumulativa, em \textbf{cada} valor de abscissa ( eixo x ), o valor apresentado no eixo das ordenadas ( eixo y ) é o valor acumulado da variável ( Conceito Enade ).
Em linhas gerais, em um gráfico desta mesma tipologia, poligonais mais à esquerda representam, no caso dos dados apresentados nesta seção, as grandes regiões com a pior distribuição, e poligonais mais à direita, as grandes regiões com os melhores conceitos.
A partir destas diretrizes, no Gráfico 4.1, é apresentada a distribuição cumulativa do Conceito Enade, segundo a grande região.
No Gráfico 4.1, observando a linha vertical pontilhada, que ajuda a visualizar o percentual de conceitos acumulados até o valor 3, ou seja, conceitos 1, 2 e 3, pode -se notar, por exemplo, que as regiões Sul ( linha rosa ) e Sudeste ( linha vermelha ) apresentam 78,5 \% dos cursos com conceito 3 ou abaixo deste valor.
A região Norte ( linha azul ), por outro lado, apresenta 100,0 \% dos cursos na mesma situação : com conceito 3 ou abaixo.
Destaca -se que, em algumas regiões, o conceito 5 não foi alcançado por nenhum curso, como ocorre nas regiões Norte ( linha azul ) e Nordeste ( linha verde ), nas quais todos os cursos têm conceito menor ou igual a 3 na região Norte e menor ou igual a 4 na região Nordeste, e as poligonais cumulativas alcançam 100 \% nos conceitos 3 e 4, respectivamente.
A região Norte ( linha azul ) apresenta os cursos com conceitos mais baixos, ao passo que a região Centro-Oeste ( linha laranja ) apresenta a distribuição com valores maiores.
Os cursos em IES situadas na região Nordeste ( linha verde ) apresentam uma situação intermediária. 4.2 CONCEITOS POR CATEGORIA ADMINISTRATIVA, MODALIDADE DE OFERTA DOS CURSOS E GRANDE REGIÃO Na Tabela 4.2, é apresentada a distribuição dos cursos de Engenharia Química participantes do Enade 2023, por categoria administrativa e modalidade de oferta, de acordo com os conceitos por eles alcançados, segundo as grandes regiões brasileiras.
De acordo com as informações da Tabela 4.2, dos 185 cursos participantes, 66 ( 35,7 \% ) eram ministrados em instituições públicas e 119 ( 64,3 \% ), em instituições privadas.
Quanto à modalidade de oferta, 182 ( 98,4 \% ) cursos eram presenciais e 3 ( 1,6 \% ), a distância.
De acordo com as informações da mesma tabela, em termos nacionais, entre instituições públicas e privadas, observa -se que, dos 9 cursos avaliados com conceito 5, 9 foram oferecidos em IES públicas, e 0, em IES privadas.
Dos 119 cursos participantes de IES privadas, o conceito 2 foi o valor modal, atribuído a 51 ( 42,9 \% ) cursos.
Em relação aos demais cursos de instituições privadas, 37 cursos ( 31,1 \% ) foram avaliados com o conceito 3, oito cursos ( 6,7 \% ) foram avaliados com o conceito 1, sete cursos ( 5,9 \% ) foram avaliados com o conceito 4 e, nenhum curso ( 0,0 \% ), com o conceito 5.
Nessa categoria, 16 cursos ( 13,4 \% ) ficaram sem conceito ( SC ).
Na rede pública, o conceito modal foi o 3, atribuído a 33 cursos, correspondentes a 50,0 \% dos 66 cursos da categoria.
Dos demais cursos participantes, 22 cursos ( 33,3 \% ) obtiveram o conceito 4, nove cursos ( 13,6 \% ) receberam o conceito 5 e dois cursos ( 3,0 \% ), o conceito 2.
O conceito 1 foi atribuído a nenhum curso ( 0,0 \% ), conforme mencionado.
Nessa categoria administrativa, nenhum curso ( 0,0 \% ) permaneceu sem conceito ( SC ).
Dos 182 cursos participantes da modalidade presencial, o conceito 3 foi o valor modal, atribuído a 69 ( 37,9 \% ) cursos.
Dos demais cursos dessa modalidade de oferta, 51 cursos ( 28,0 \% ) receberam o conceito 2 ; 29 cursos ( 15,9 \% ) receberam o conceito 4 ; nove cursos ( 4,9 \% ), o conceito 5 ; e oito cursos ( 4,4 \% ), o conceito 1.
Nessa modalidade de oferta, 16 cursos ( 8,8 \% ) ficaram sem conceito ( SC ).
Dos 3 cursos participantes da modalidade a distância, o conceito 2 foi o valor modal, atribuído a 2 ( 66,7 \% ) cursos.
Na análise por região, observa -se, que, na região Norte, as instituições públicas participaram com cinco cursos ( 71,4 \% total regional ), dos quais nenhum ficou sem conceito ( SC ) ou recebeu o conceito 5.
O conceito modal foi 3, atribuído a cinco cursos, correspondendo a 100,0 \% dessa categoria na região.
As instituições privadas participaram com 2 cursos na região Norte ( 28,6 \% do total regional ), e nenhum curso ficou sem conceito ( SC ).
O conceito modal nas instituições privadas na região foi o 1, atribuído a um curso, o que corresponde a 50,0 \% dessa categoria na região.
Os conceitos 1, 5 e 4 foram atribuídos a um, nenhum e nenhum curso, respectivamente.
O conceito 3 foi atribuído a nenhum curso ( 0,0 \% ).
Nenhum curso ficou sem conceito ( SC ) ou recebeu os demais conceitos.
Dos cursos oferecidos na região Norte, sete eram presenciais e nenhum ficou sem conceito ( SC ).
Dos sete cursos oferecidos, cinco obtiveram o conceito 3 ( modal ) ; um, o conceito 2 e um, o conceito 1 Nenhum curso era na modalidade a distância.
Na região Nordeste, a rede privada concentra 17 dos 29 cursos participantes, o equivalente a 58,6 \% do total da região.
Entre os cursos oferecidos por IES privadas no Nordeste, sete cursos foram avaliados com conceito 2, o conceito modal ; dois cursos obtiveram o conceito 3 ; um curso, o conceito 4 ; um curso, o conceito 1 ; e nenhum curso obteve o conceito 5.
Nessa categoria administrativa e grande região, seis cursos ficaram sem conceito ( SC ).
As instituições públicas da região Nordeste participaram com 12 cursos ( 41,4 \% do total da região ).
Desses, sete ( 58,3 \% ) obtiveram conceito 3, o valor modal, seguidos por cinco cursos, que receberam o conceito 4.
O conceito 5 foi atribuído a nenhum curso ; o conceito 2, a nenhum curso ; e o conceito 1, a nenhum curso.
Nenhum curso ficou sem conceito ( SC ).
Dos cursos oferecidos na região Nordeste, todos ( 29 ) eram presenciais : seis ficaram sem conceito ( SC ), um obteve o conceito 1 ; sete, o conceito 2 ; nove, o conceito 3, o valor modal ; seis, o conceito 4 ; e nenhum, o conceito 5.
Nenhum curso era na modalidade a distância.
Na região Sudeste, a proporção de cursos da rede privada, 71,3 \%, correspondeu a 62 dos 87 cursos participantes.
Nessa combinação de categoria administrativa e grande região, o conceito modal foi o 2, atribuído a 28 cursos ; e oito cursos ficaram sem conceito ( SC ).
Os demais foram avaliados com os seguintes conceitos : 3 ( 21 cursos ), 1 ( três cursos ) e 4 ( dois cursos ).
Dos 25 cursos oferecidos por instituições públicas na região Sudeste, quatro obtiveram o conceito 5 ; 11, o conceito 4 ; dez, o conceito 3 ; nenhum obteve conceito 2 e nenhum obteve o conceito 1.
Nenhum curso ficou sem conceito ( SC ).
Ainda na região Sudeste, 100 \% dos cursos foram oferecidos na modalidade presencial, 87.
Para esses, a moda foi o conceito 3, com 31 cursos.
Os demais cursos presenciais da região foram alocados nos conceitos 2 ( 28 cursos ), 4 ( 13 cursos ), 5 ( quatro cursos ) e 1 ( três cursos ).
Oito cursos presenciais ficaram sem conceito ( SC ).
Na região Sudeste, nenhum curso foi oferecido na modalidade a distância.
As instituições privadas concentraram 36 dos 56 cursos participantes da região Sul, 64,3 \% do total regional.
Nessa combinação de categoria administrativa e grande região, 14 receberam o conceito 2 e 14 receberam o conceito 3 ( bimodal ) ; quatro receberam o conceito 4 ; dois, o conceito 1 ; e nenhum obteve o conceito 5.
Dois cursos ficaram sem conceito ( SC ).
As instituições públicas na região Sul participaram com 20 cursos ( 35,7 \% dos cursos da região ), aos quais foram atribuídos os conceitos 2 ( dois cursos ), 3 ( dez cursos ), 4 ( quatro cursos ) e 5 ( quatro cursos ).
Nenhum curso da região Sul oferecido por IES públicas ficou sem conceito ( SC ).
Na região Sul, 53 dos cursos foram oferecidos na modalidade presencial.
Para esses, a moda foi o conceito 3, com 23 cursos.
Os demais cursos presenciais da região foram alocados nos conceitos 2 ( 14 cursos ), 4 ( oito cursos ), 5 ( quatro cursos ) e 1 ( dois cursos ).
Ficaram sem conceito ( SC ) dois cursos.
Na Na modalidade a distância, foram oferecidos três cursos, sendo dois cursos com conceito 2 e um curso com conceito 3.
Na região Centro-Oeste, 2 dos 6 cursos participantes eram de IES privadas ( 33,3 \% em termos regionais ).
Desses, um curso recebeu o conceito 1 e outro curso recebeu o conceito 2.
Quanto aos quatro cursos oferecidos por instituições públicas na região Centro-Oeste ( 66,7 \% do total regional ), um ficou com conceito 3 ; dois, com conceito 4 ; e um ficou com conceito 5.
Nenhum ficou sem conceito ( SC ) nessa região.
Também na região Centro-Oeste, todos os seis cursos foram ofertados na modalidade presencial.
Para esses, a moda foi o conceito 4, com 2 cursos ( 33,3 \% ).
Os demais cursos presenciais da região foram alocados nos conceitos 5 ( um curso, 16,7 \% ), 3 ( um curso, 16,7 \% ), 2 ( um curso, 16,7 \% ) e 1 ( um curso, 16,7 \% ).
Nenhum curso presencial da região ficou sem conceito ( SC ).
Não houve curso oferecido na modalidade a distância.
A título de síntese, no Gráfico 4.2, apresenta -se a distribuição do Conceito Enade, segundo a categoria administrativa da IES.
Nesse gráfico, poligonais mais à esquerda representam Conceito Enade mais baixos, e poligonais mais à direita, Conceito Enade mais altos.
É importante observar que, para \textbf{cada} poligonal presente nesse gráfico, a porcentagem constante no eixo y refere -se ao universo de cursos oferecidos por instituições públicas ( poligonal azul ) ou privadas ( poligonal verde ), respectivamente.
Os cursos em IES públicas ( poligonal azul ) apresentam uma poligonal mais à direita do que os em IES privadas ( poligonal verde ) e, consequentemente, uma distribuição de Conceito Enade com valores maiores. 4.3 CONCEITOS POR ORGANIZAÇÃO ACADÊMICA E GRANDE REGIÃO Na Tabela 4.3, encontra -se a distribuição dos conceitos atribuídos aos cursos participantes do Enade 2023, na área de Engenharia Química, por organização acadêmica, segundo as grandes regiões brasileiras.
Como mostram os dados da Tabela 4.3, dos 185 cursos de Engenharia Química participantes, 110 eram oferecidos em universidades, 41 em centros universitários e 29 em faculdades ; cinco cursos foram oferecidos por CEFET/IF.
Essa distribuição corresponde a, respectivamente, 59,4 \%, 22,2 \%, 15,7 \% e 2,7 \% dos cursos.
De acordo com os dados apresentados na mesma tabela, de todos os 29 cursos avaliados com o conceito 4, 25 eram vinculados a universidades.
Os cursos vinculados a esse tipo de organização acadêmica tiveram o conceito 3 como conceito modal, atribuído a 50 cursos ( 45,5 \% ).
Os demais cursos ofertados por universidades receberam os conceitos 2 ( 22 cursos ), 5 ( oito cursos ) e 1 ( um curso ).
Quatro cursos ficaram sem conceito ( SC ).
Dos cursos vinculados a centros universitários, o conceito modal foi o 2, atribuído a 20 cursos ( 48,8 \% ).
Seis cursos vinculados a esse tipo de organização acadêmica ficaram sem conceito ( SC ), e os outros cursos receberam os conceitos 3 ( onze cursos ), 1 ( três cursos ), 4 ( um curso ) e 5 ( nenhum curso ).
Dos 29 cursos mantidos por faculdades, 11 ( 37,9 \% ) receberam conceito 2, o conceito modal.
O segundo conceito mais frequentemente atribuído aos cursos oferecidos por essa categoria de IES foi o conceito 3, recebido por sete cursos.
Na sequência, quatro cursos obtiveram conceito 1 ; um curso, o conceito 5 ; e nenhum curso, o conceito 4.
Dos cursos oferecidos em faculdades, seis cursos \textbf{estiveram} sem conceito ( SC ).
Dos cursos oferecidos por CEFET/IF, o conceito modal foi o 4, atribuído a três cursos ( 60,0 \% ) e dois cursos receberam o o conceito 3.
Considerando -se, separadamente, as regiões brasileiras, verifica -se que, na região Norte, as universidades participaram com cinco cursos avaliados, todos alcançaram o conceito 3.
Os centros universitários da região Norte foram representados por dois cursos, dos quais um curso obteve conceito 1 e um curso, conceito 2.
Não houve cursos participantes oferecidos por faculdades e CEFET/IF Na região Nordeste, as universidades participaram com 15 dos 29 cursos da área de Engenharia Química oferecidos.
O conceito modal foi 3, atribuído a oito cursos.
Outros quatro cursos obtiveram o conceito 4 ; três cursos, o conceito 2 ; nenhum curso, o conceito 5 ; e nenhum curso, o conceito 1.
Nenhum curso ficou sem conceito ( SC ).
Os centros universitários contaram com oito cursos participantes na região Nordeste, três dos quais receberam o conceito modal 2.
Para os demais cursos, foram atribuídos os conceitos 4 ( um curso ) e 1 ( um curso).Três cursos oferecidos por centros universitários no Nordeste ficaram sem conceito ( SC ).
As faculdades foram representadas por cinco cursos na região Nordeste, os quais receberam os 2 ( um curso ), 3 ( um curso ) e três cursos ficaram sem conceito ( SC ).
Na região Nordeste, os CEFET/IF participaram com um curso, ao qual foi atribuído o conceito 4.
Na região Sudeste, as universidades concentraram 44 dos 87 cursos de Engenharia Química da região.
Dos cursos oferecidos em IES com este tipo de organização acadêmica, nessa região, o conceito modal foi 3, atribuído a 16 cursos.
Os demais cursos receberam os conceitos 4 ( doze cursos ), 2 ( dez cursos ), 5 ( três cursos ) e três cursos ficaram sem conceito ( SC ).
Os centros universitários participaram com 22 cursos na região Sudeste, dos quais dez obtiveram o conceito modal 3.
O conceito 2 foi atribuído a nove cursos ; o conceito 1 a um curso e dois cursos ficaram sem conceito ( SC ) nessa combinação de organização acadêmica e região.
Já as faculdades foram representadas por 19 cursos na região Sudeste, dos quais três ficaram sem conceito ( SC ).
O conceito modal 2 foi atribuído a 9 cursos.
Aos demais cursos foram atribuídos os conceitos 3 ( quatro cursos ), 1 ( dois cursos ) e 5 ( um curso ).
Os CEFET/IF participaram com dois cursos na região Sudeste, aos quais foram atribuídos os conceitos 3 e 4.
Dos 56 cursos da região Sul, 42 eram mantidos por universidades, para os quais o conceito modal foi 3, atribuído a 20 cursos.
Os demais cursos receberam os conceitos 2 ( nove cursos ), 4 ( sete cursos ), 5 ( quatro cursos ) e 1 ( um curso ).
Em IES com esse tipo de organização acadêmica, um curso da região Sul ficou sem conceito ( SC ).
Dos nove cursos integrantes de centros universitários da região Sul, um ficou sem conceito ( SC ), e sete cursos receberam conceito 2 e um curso recebeu o conceito 3.
Quanto aos três cursos vinculados a faculdades na região Sul, dois cursos receberam o conceito 3, valor modal e um curso recebeu o conceito 1.
Os CEFET/IF participaram com dois cursos na região Sul, aos quais foram atribuídos os seguintes conceitos 3 e 4.
Na região Centro-Oeste, 4 dos 6 cursos oferecidos eram mantidos por universidades.
Para esse tipo de organização acadêmica, o conceito modal foi 4, atribuído a dois cursos.
Um curso recebeu o conceito 5 e um curso recebeu o conceito 3.
Os centros universitários da região Centro-Oeste contaram com nenhum curso avaliado.
Dos dois cursos oferecidos em faculdades na região Centro-Oeste, um curso recebeu o conceito 2 e um curso recebeu o conceito 1.
No que se refere aos CEFET/IF da região Centro-Oeste, nenhum curso foi avaliado.
De maneira semelhante à forma como se procedeu em relação à seção anterior, esta seção encerra -se com o Gráfico 4.3, que apresenta a distribuição do Conceito Enade dos cursos de Engenharia Química segundo a organização acadêmica da IES.
No Gráfico 4.3, apresenta -se a distribuição do Conceito Enade segundo a organização acadêmica da IES.
A leitura desse último gráfico é semelhante aos anteriores, seguindo -se a lógica de que poligonais deslocadas para a esquerda correspondem aos conceitos mais baixos ( 1, 2 e 3 ) e poligonais deslocadas para a direita correspondem a Conceitos Enade mais altos.
Nas universidades ( linha azul ), todos os cursos obtiveram conceito 5 ou abaixo disso, já que a poligonal cumulativa alcança 100 \% para o conceito 5.
Além disso, os cursos de universidades apresentam uma poligonal mais à direita do que os Centro Universitários e faculdades, denotando uma melhor distribuição de Conceitos Enade.
No outro extremo, encontram -se os cursos mantidos por CEFET/IF ( linha laranja ), apresentando as respectivas poligonais mais à direita e, consequentemente, uma distribuição de Conceitos Enade com os maiores valores.

% matched lemmas: cada, estar
\end{textsample}
